% GENERATED FILE. Edit canonical lessons, then run npm run generate:books.
% canonical-source-hash: fnv1a64:7ae9f61053294217
% canonical-lessons: HI-C370-sabha, HI-C370-pradhan, HI-C370-zila, HI-C370-parti, HI-C370-taykarna

\chapter{The Village Meeting}
\label{ch:hi-the-village-meeting}
\hlchaptermodality{370}

\begin{chapteropening}
\textbf{By the end of this chapter:} \emph{I can say who decided what at a meeting.}

Complete the last lesson of chapter 370: I can say who decided what at a meeting.
\end{chapteropening}

\section[\texorpdfstring{\hi{सभा} (sabhā)}{सभा (sabhā)}]{\hi{सभा} (sabhā) — an assembly, a public meeting}

\label{lesson:HI-C370-sabha}



\begin{quote}
Before the new one: say the Hindi for an employee, then the Hindi for a ban.
\end{quote}

\subsection*{You'll want to know: \hi{सभा}}
\textbf{\hi{सभा}} — \emph{sabhā} — \textquotedblleft{}an assembly, a public meeting\textquotedblright{}.

\begin{quote}
\hi{शाम} \hi{को} \hi{गाँव} \hi{में} \hi{एक} \hi{सभा} \hi{होगी।}
\end{quote}

\begin{quote}
\emph{śām ko gā̃v meṁ ek sabhā hogī.}
\end{quote}

\begin{quote}
There will be a public meeting in the village this evening.
\end{quote}

Feminine.

\begin{cousinweb}
\textbf{\hi{सभा}} (\emph{sabhā}) is a tatsama, borrowed from Sanskrit \emph{sabhā}, \textquotedblleft{}an assembly\textquotedblright{} (Monier-Williams).
\end{cousinweb}

\begin{grammarlens}[title={What you've built}]
One more word for giving reasons and opinions.
\end{grammarlens}

\subsection*{Your turn}
\begin{itemize}
\raggedright
  \item \emph{Say it:} \emph{sabhā}
  \item \emph{Say it:} \emph{sabhā}, once more
  \item \emph{Say it:} \emph{śām ko gā̃v meṁ ek sabhā hogī.} — There will be a public meeting in the village this evening.
  \item \emph{Recall:} say the Hindi for an employee, then the Hindi for a ban, then say \emph{sabhā} again
\end{itemize}

\begin{tcolorbox}[breakable,colback=teal!4,colframe=teal!35!black,title={Before you move on}]
What does \hi{सभा} mean? (\textquotedblleft{}An assembly, a public meeting\textquotedblright{}.) Say it once more.
\end{tcolorbox}
\section[\texorpdfstring{\hi{प्रधान} (pradhān)}{प्रधान (pradhān)}]{\hi{प्रधान} (pradhān) — a village head; a head, a chief}

\label{lesson:HI-C370-pradhan}



\begin{quote}
Before the new one: say the Hindi for a ban, then the Hindi for an assembly.
\end{quote}

\subsection*{You'll want to know: \hi{प्रधान}}
\textbf{\hi{प्रधान}} — \emph{pradhān} — \textquotedblleft{}a village head; a head, a chief\textquotedblright{}.

\begin{quote}
\hi{गाँव} \hi{के} \hi{प्रधान} \hi{ने} \hi{सभा} \hi{बुलाई।}
\end{quote}

\begin{quote}
\emph{gā̃v ke pradhān ne sabhā bulāī.}
\end{quote}

\begin{quote}
The village head called a meeting.
\end{quote}

Masculine or feminine, by the person meant.

\begin{cousinweb}
\textbf{\hi{प्रधान}} (\emph{pradhān}) is a tatsama, borrowed from Sanskrit \emph{pradhāna}, \textquotedblleft{}the chief thing or person\textquotedblright{} (Monier-Williams).
\end{cousinweb}

\begin{grammarlens}[title={What you've built}]
One more word for giving reasons and opinions.
\end{grammarlens}

\subsection*{Your turn}
\begin{itemize}
\raggedright
  \item \emph{Say it:} \emph{pradhān}
  \item \emph{Say it:} \emph{pradhān}, once more
  \item \emph{Say it:} \emph{gā̃v ke pradhān ne sabhā bulāī.} — The village head called a meeting.
  \item \emph{Recall:} say the Hindi for a ban, then the Hindi for an assembly, then say \emph{pradhān} again
\end{itemize}

\begin{tcolorbox}[breakable,colback=teal!4,colframe=teal!35!black,title={Before you move on}]
What does \hi{प्रधान} mean? (\textquotedblleft{}A village head; a head, a chief\textquotedblright{}.) Say it once more.
\end{tcolorbox}
\section[\texorpdfstring{\hi{ज़िला} (zilā)}{ज़िला (zilā)}]{\hi{ज़िला} (zilā) — a district}

\label{lesson:HI-C370-zila}



\begin{quote}
Before the new one: say the Hindi for an assembly, then the Hindi for a village head.
\end{quote}

\subsection*{You'll want to know: \hi{ज़िला}}
\textbf{\hi{ज़िला}} — \emph{zilā} — \textquotedblleft{}a district\textquotedblright{}.

\begin{quote}
\hi{हमारा} \hi{ज़िला} \hi{पहाड़ों} \hi{में} \hi{है।}
\end{quote}

\begin{quote}
\emph{hamārā zilā pahāṛoṁ meṁ hai.}
\end{quote}

\begin{quote}
Our district is in the mountains.
\end{quote}

Masculine.

\begin{cousinweb}
\textbf{\hi{ज़िला}} (\emph{zilā}) is usually traced to Arabic \emph{ḍilʿ}, \textquotedblleft{}a side\textquotedblright{}, which reached Hindi through Persian. Persian and Hindi use it for a district.
\end{cousinweb}

\begin{grammarlens}[title={What you've built}]
One more word for giving reasons and opinions.
\end{grammarlens}

\subsection*{Your turn}
\begin{itemize}
\raggedright
  \item \emph{Say it:} \emph{zilā}
  \item \emph{Say it:} \emph{zilā}, once more
  \item \emph{Say it:} \emph{hamārā zilā pahāṛoṁ meṁ hai.} — Our district is in the mountains.
  \item \emph{Recall:} say the Hindi for an assembly, then the Hindi for a village head, then say \emph{zilā} again
\end{itemize}

\begin{tcolorbox}[breakable,colback=teal!4,colframe=teal!35!black,title={Before you move on}]
What does \hi{ज़िला} mean? (\textquotedblleft{}A district\textquotedblright{}.) Say it once more.
\end{tcolorbox}
\section[\texorpdfstring{\hi{पार्टी} (pārṭī)}{पार्टी (pārṭī)}]{\hi{पार्टी} (pārṭī) — a political party; a party}

\label{lesson:HI-C370-parti}



\begin{quote}
Before the new one: say the Hindi for a village head, then the Hindi for a district.
\end{quote}

\subsection*{You'll want to know: \hi{पार्टी}}
\textbf{\hi{पार्टी}} — \emph{pārṭī} — \textquotedblleft{}a political party; a party\textquotedblright{}.

\begin{quote}
\hi{उसकी} \hi{पार्टी} \hi{चुनाव} \hi{जीत} \hi{गई।}
\end{quote}

\begin{quote}
\emph{uskī pārṭī cunāv jīt gaī.}
\end{quote}

\begin{quote}
His party won the election.
\end{quote}

Feminine.

\begin{cousinweb}
\textbf{\hi{पार्टी}} (\emph{pārṭī}) is a loan from English \emph{party}.
\end{cousinweb}

\begin{grammarlens}[title={What you've built}]
One more word for giving reasons and opinions.
\end{grammarlens}

\subsection*{Your turn}
\begin{itemize}
\raggedright
  \item \emph{Say it:} \emph{pārṭī}
  \item \emph{Say it:} \emph{pārṭī}, once more
  \item \emph{Say it:} \emph{uskī pārṭī cunāv jīt gaī.} — His party won the election.
  \item \emph{Recall:} say the Hindi for a village head, then the Hindi for a district, then say \emph{pārṭī} again
\end{itemize}

\begin{tcolorbox}[breakable,colback=teal!4,colframe=teal!35!black,title={Before you move on}]
What does \hi{पार्टी} mean? (\textquotedblleft{}A political party; a party\textquotedblright{}.) Say it once more.
\end{tcolorbox}
\section[\texorpdfstring{\hi{तय} \hi{करना} (tay karnā)}{तय करना (tay karnā)}]{\hi{तय} \hi{करना} (tay karnā) — to fix (a date), to decide}

\label{lesson:HI-C370-taykarna}



\begin{quote}
Before the new one: say the Hindi for a district, then the Hindi for a political party.
\end{quote}

\subsection*{You'll want to know: \hi{तय} \hi{करना}}
\textbf{\hi{तय} \hi{करना}} — \emph{tay karnā} — \textquotedblleft{}to fix (a date), to decide\textquotedblright{}.

\begin{quote}
\hi{सभा} \hi{ने} \hi{अगली} \hi{बैठक} \hi{की} \hi{तारीख़} \hi{तय} \hi{की।}
\end{quote}

\begin{quote}
\emph{sabhā ne aglī baiṭhak kī tārīkh tay kī.}
\end{quote}

\begin{quote}
The meeting settled the date of the next one.
\end{quote}

Past, with \hi{ने} (he or she …): \hi{उसने} \hi{दिन} \hi{तय} \hi{किया}, \emph{usne din tay kiyā}.

\begin{cousinweb}
\textbf{\hi{तय} \hi{करना}} (\emph{tay karnā}) joins \hi{तय} to \hi{करना}, \textquotedblleft{}to do\textquotedblright{}. \hi{तय} is usually traced to Arabic \emph{ṭayy}, \textquotedblleft{}folding up\textquotedblright{}, which reached Hindi through Persian. The sense of settling a matter is a Persian and Hindi one.
\end{cousinweb}

\begin{grammarlens}[title={What you've built}]
One more word for giving reasons and opinions.
\end{grammarlens}

\subsection*{Your turn}
\begin{itemize}
\raggedright
  \item \emph{Say it:} \emph{tay karnā}
  \item \emph{Say it:} \emph{tay karnā}, once more
  \item \emph{Say it:} \emph{sabhā ne aglī baiṭhak kī tārīkh tay kī.} — The meeting settled the date of the next one.
  \item \emph{Recall:} say the Hindi for a district, then the Hindi for a political party, then say \emph{tay karnā} again
\end{itemize}

\begin{tcolorbox}[breakable,colback=teal!4,colframe=teal!35!black,title={Before you move on}]
What does \hi{तय} \hi{करना} mean? (\textquotedblleft{}To fix (a date), to decide\textquotedblright{}.) Say it once more.
\end{tcolorbox}
